\chapter{Implementación y resultados}
\label{ch:implementacion}

Este capítulo describe el proceso de implementación de cada componente del simulador: las decisiones técnicas concretas tomadas durante el desarrollo, los problemas encontrados y cómo se resolvieron, y los resultados obtenidos en el entrenamiento y evaluación de los modelos de elección modal. Es el capítulo más extenso de la memoria porque concentra el trabajo de ingeniería real que da soporte a la arquitectura descrita en el capítulo anterior.

% -------------------------------------------------------
\section{Infraestructura de enrutado viario: OSRM}
\label{sec:impl_osrm}
% -------------------------------------------------------

\subsection{Origen y características del extracto OSM}

La fuente de datos viarios del simulador es el extracto de OpenStreetMap para Castilla-La Mancha, distribuido por Geofabrik en formato \texttt{.osm.pbf} \cite{GeofabrikDownloads}. Este formato binario comprimido es el estándar de distribución de extractos OSM y es directamente compatible con el pipeline de preprocesado de OSRM. El archivo descargado tiene un tamaño aproximado de 98 MB, lo que lo hace manejable para almacenamiento local y preprocesado en máquina de desarrollo, a diferencia del extracto de España (1,3 GB) o el de Europa completo (32 GB).

La elección de Castilla-La Mancha está motivada tanto por el ámbito académico del TFM (Escuela Superior de Informática de Albacete, UCLM) como por la disponibilidad del feed GTFS de buses urbanos de Toledo, que es la ciudad en la que se concentra el caso de estudio. Utilizar el extracto regional en lugar del nacional supone una reducción notable del tiempo de preprocesado y del espacio en disco de los grafos generados, sin penalizar la calidad del enrutado dentro del área urbana de Toledo.

\subsection{Evolución del despliegue: servidor remoto a infraestructura local}

La integración de OSRM en el simulador pasó por tres fases con complejidad creciente, que reflejan las limitaciones encontradas y las decisiones tomadas para superarlas.

En la primera fase se usó el servidor de demostración oficial de OSRM (\texttt{router.project-osrm.org}). Este servidor acepta peticiones del tipo \texttt{/route/v1/\{profile\}/lon1,lat1;lon2,lat2} y es adecuado para pruebas iniciales. Sin embargo, aunque la API permite especificar el perfil en la URL, el servidor de demostración está capado y devuelve en la práctica el mismo resultado para coche, bicicleta y a pie. Dado que el objetivo del simulador es precisamente comparar modos de transporte con atributos diferenciados, esta limitación hacía inviable el uso de este servidor.

En la segunda fase se recurrió a las instancias públicas de FOSSGIS, que sí ofrecen perfiles diferenciados (\texttt{routing.openstreetmap.de/routed-car}, \texttt{.../routed-bike}, \texttt{.../routed-foot}). Estas instancias permitieron verificar la viabilidad del enfoque multi-perfil antes de comprometerse con un despliegue local. No obstante, los servicios públicos introducen latencia variable, dependencia de disponibilidad externa y no permiten fijar una versión de los datos, lo que dificulta la reproducibilidad.

En la tercera y definitiva fase se desplegó OSRM localmente como contenedores Docker, uno por perfil. Esta decisión sigue la misma lógica que motivó el uso de OSRM frente a APIs comerciales (reproducibilidad, control de versión, sin coste por petición), pero aplicada ahora también a la gestión de perfiles. La imagen oficial de OSRM (\texttt{osrm/osrm-backend}) incluye los perfiles estándar \texttt{car.lua}, \texttt{bicycle.lua} y \texttt{foot.lua}, lo que elimina la necesidad de obtenerlos o configurarlos externamente.

\subsection{Pipeline de preprocesado por perfil}

El preprocesado OSRM consta de tres etapas que deben ejecutarse una sola vez por perfil, o cada vez que cambie el extracto OSM. Todas las etapas se ejecutan con la imagen oficial de Docker para garantizar versión fija del software:

\begin{enumerate}
  \item \textbf{Extracción} (\texttt{osrm-extract}): transforma el \texttt{.osm.pbf} en una representación interna del grafo de red viaria, aplicando las restricciones de velocidad, giros y acceso definidas en el archivo de perfil. Para el perfil de coche se usa \texttt{/opt/car.lua}, para bicicleta \texttt{/opt/bicycle.lua} y para peatón \texttt{/opt/foot.lua}.
  \item \textbf{Particionado} (\texttt{osrm-partition}): divide el grafo en celdas jerárquicas para el algoritmo MLD. Esta etapa es específica del algoritmo Multi-Level Dijkstra y es necesaria antes de customizar.
  \item \textbf{Customización} (\texttt{osrm-customize}): ajusta los pesos de las aristas en función del perfil sobre la estructura de celdas generada en la etapa anterior.
\end{enumerate}

El resultado de las tres etapas es el conjunto de archivos \texttt{clm.osrm} y sus ficheros asociados, que se montan como volumen Docker en el contenedor de explotación. Cada perfil mantiene su propio directorio (\texttt{osrm-clm/car/}, \texttt{osrm-clm/bike/}, \texttt{osrm-clm/foot/}), por lo que los tres conjuntos de artefactos coexisten sin interferencia.

\subsection{Despliegue y verificación}

El servidor de enrutado se lanza con \texttt{osrm-routed --algorithm mld /data/clm.osrm}. Los tres contenedores exponen el puerto 5000 internamente pero se mapean a puertos distintos del anfitrión (5000, 5001 y 5002), lo que permite que el backend FastAPI los distinga por URL. En producción (Docker Compose), los contenedores se referencian por nombre de servicio interno (\texttt{osrm-car:5000}, \texttt{osrm-bike:5000}, \texttt{osrm-foot:5000}), por lo que el cambio de nombre de host es transparente para el backend.

La verificación del servicio se realiza mediante una petición directa al endpoint de ruta con coordenadas conocidas dentro del área de Toledo. Una respuesta válida incluye la distancia en metros, la duración en segundos y la geometría codificada en polyline. Si el preprocesado se hizo con un perfil incorrecto o el grafo está incompleto, OSRM devuelve código de estado \texttt{NoRoute} o una ruta degradada fácilmente detectable.

% -------------------------------------------------------
\section{Planificación multimodal: OpenTripPlanner}
\label{sec:impl_otp}
% -------------------------------------------------------

\subsection{Fuente de datos GTFS y proceso de selección}

La integración del transporte público requirió localizar un feed GTFS adecuado para el caso de estudio. El proceso de selección es parte del trabajo de implementación porque no existe un único repositorio de datos de transporte público en España con cobertura y calidad garantizadas para todos los operadores.

El canal de búsqueda fue el Punto de Acceso Nacional de Transporte Multimodal (NAP), gestionado por el Ministerio de Transportes y Movilidad Sostenible, que agrega información de oferta de transporte de viajeros de múltiples operadores en línea con el marco europeo de datos de movilidad \cite{NAPHome}. El NAP se consolidó como la fuente de referencia a partir de su presentación pública en junio de 2021.

Durante el desarrollo se exploraron varios feeds antes de decidir el feed definitivo:

\begin{itemize}
  \item \textbf{EMT Madrid} \cite{NAPEMTMadrid}: feed de gran volumen (4.914 paradas, 236 rutas, 87.137 viajes). Permitió validar la integración del backend GTFS y del planificador OTP con datos reales, pero el tamaño implica usar el extracto de España completo para OSM, lo que aumenta el coste de preprocesado. Se usó como banco de pruebas previo al caso definitivo.
  \item \textbf{EMT Valencia} \cite{NAPEMTValencia}: consultado brevemente como referencia alternativa, descartado en favor de una ciudad en Castilla-La Mancha.
  \item \textbf{Buses urbanos de Toledo (UNAUTO S.L.)} \cite{NAPToledoUrbano}: feed definitivo. 268 paradas, 56 rutas y 46.597 viajes. Incluye geometría de rutas (shapes), información de accesibilidad y cobertura del casco urbano de Toledo, el polígono industrial y urbanizaciones periféricas. El operador era inicialmente «Autobuses Toledo» en el momento de la primera descarga; durante el desarrollo cambió a UNAUTO S.L. (Grupo Ruiz), aunque el formato GTFS permaneció idéntico (\texttt{agency.txt}, \texttt{calendar\_dates.txt}, \texttt{routes.txt}, \texttt{shapes.txt}, \texttt{stop\_times.txt}, \texttt{stops.txt}, \texttt{trips.txt}).
\end{itemize}

Toledo fue la elección final por su tamaño manejable, cobertura del área de interés del TFM y disponibilidad en Castilla-La Mancha, lo que permite usar el mismo extracto OSM regional que OSRM.

\subsection{Vigencia del feed y decisión de fecha fija}

Los feeds GTFS del NAP tienen vigencia limitada: para el caso de Toledo, los archivos cubren aproximadamente mes y medio de servicio programado. Esto genera un problema operativo en el simulador: si la fecha de consulta cae fuera del rango de validez del feed cargado en OTP, el planificador no devuelve itinerarios reales y el resultado es inservible para la inferencia modal.

El problema se manifestó al volver de las vacaciones de Navidades de 2025: la fecha del sistema había superado el rango cubierto por el feed descargado, por lo que OTP dejó de devolver itinerarios. La solución adoptada fue establecer una fecha fija en el backend (\texttt{date=2025-12-01}, \texttt{time=12:00}), que queda siempre dentro del rango del feed en uso, y actualizar el feed junto con la fecha si fuera necesario en el futuro. Esta decisión garantiza reproducibilidad: todos los experimentos y ejecuciones del simulador trabajan sobre el mismo escenario de oferta de transporte.

La hora de consulta se fijó a las 12:00 para evitar itinerarios nocturnos con servicio degradado o inexistente, que distorsionarían el vector de variables de ruta de forma no representativa.

\subsection{Construcción del grafo multimodal}

El grafo multimodal de OTP se genera a partir del comando de construcción (\texttt{--build --save}) sobre un directorio que contiene el extracto OSM de Castilla-La Mancha y el archivo ZIP del feed GTFS de Toledo. OTP 2.5.0 combina ambas fuentes en un único grafo (\texttt{graph.obj}) que integra la red viaria para los tramos de acceso a pie y la red programada del autobús para los tramos en vehículo.

El proceso de construcción tiene una duración aproximada de dos minutos en la máquina de desarrollo. Al igual que los grafos OSRM, el \texttt{graph.obj} resultante no se versiona en el repositorio Git por su tamaño; debe generarse o descargarse antes del primer despliegue, según se detalla en el anexo~\ref{anx:despliegue}.

Una vez construido el grafo, OTP se lanza en modo de explotación (\texttt{--load --serve}) y queda disponible en el puerto 8080, exponiendo la API REST de planificación \cite{OTPRouteRequest}.

\subsection{Integración de itinerarios en el backend}

El router \texttt{/api/otp} del backend consulta OTP mediante peticiones HTTP a \texttt{/otp/routers/default/plan} y deserializa la respuesta para extraer los \textit{legs} del itinerario seleccionado. La selección del itinerario sigue este criterio: los itinerarios devueltos se ordenan por duración total; se selecciona automáticamente el primero que incluye al menos un tramo de transporte público (\texttt{transitLeg~=~true} o modo distinto de WALK, BICYCLE, CAR); si ninguno cumple la condición, se devuelve el de menor duración.

El frontend permite al usuario paginar manualmente entre itinerarios alternativos mediante los botones «Anterior» y «Siguiente», que envían el índice de itinerario deseado (\texttt{itinerary\_index}) al backend en la siguiente petición. El mismo índice se propaga al pipeline de inferencia para garantizar que el vector de características se calcula sobre el mismo itinerario que se está visualizando.

% -------------------------------------------------------
\section{Integración del feed GTFS}
\label{sec:impl_gtfs}
% -------------------------------------------------------

El router \texttt{/api/gtfs} expone la información estática de la red de autobús de Toledo directamente desde el feed descomprimido en el directorio de datos del backend. Esta capa es independiente de OpenTripPlanner: lee los archivos \texttt{stops.txt}, \texttt{routes.txt}, \texttt{trips.txt}, \texttt{stop\_times.txt}, \texttt{shapes.txt} y \texttt{calendar\_dates.txt} con pandas en el arranque del servicio y mantiene los DataFrames en memoria.

Los endpoints disponibles cubren cuatro operaciones:

\begin{itemize}
  \item \textbf{Listado de paradas} con filtro de bounding box: devuelve todas las paradas del feed, con su identificador, nombre, coordenadas y las líneas que pasan por cada una. El límite de paradas devueltas es configurable (defecto 5.000), suficiente para el feed de Toledo (268 paradas).
  \item \textbf{Listado de rutas}: devuelve el catálogo completo de líneas con identificador, nombre corto, nombre largo y atributos opcionales de color.
  \item \textbf{Detalle de ruta}: dado el identificador de una ruta, devuelve la secuencia ordenada de paradas y la geometría del shape asociado para su representación cartográfica.
  \item \textbf{Horarios por fecha}: dado un identificador de ruta y una fecha, devuelve los servicios programados agrupados por dirección, con el número de expediciones, la primera y última salida y el listado de horarios de salida.
\end{itemize}

La separación entre esta capa y OTP permite consultar la red aunque el planificador no esté disponible, lo que es útil durante el desarrollo y facilita la depuración.

La asignación de color a cada línea en la visualización cartográfica se realiza de forma determinista mediante una función de hash sobre el identificador de ruta:

\begin{equation}
  \text{color}(\mathit{route\_id}) = \text{PALETTE}\left[\left|\sum_{i} \mathit{route\_id}[i] \cdot 31^{\,\text{pos}(i)}\right| \bmod |\text{PALETTE}|\right]
\end{equation}

donde PALETTE es un vector de siete colores predefinidos. Esta estrategia garantiza que la misma línea muestra siempre el mismo color independientemente del orden de carga o del estado de la aplicación.

% -------------------------------------------------------
\section{Implementación del backend FastAPI}
\label{sec:impl_backend}
% -------------------------------------------------------

\subsection{Estructura modular}

El backend está implementado con FastAPI sobre Python y organizado en módulos por responsabilidad: un directorio \texttt{api/} con los cuatro routers (\texttt{routes\_osrm.py}, \texttt{routes\_otp.py}, \texttt{routes\_gtfs.py}, \texttt{routes\_lpmc.py}), un directorio \texttt{services/} con la lógica de negocio desacoplada de HTTP (\texttt{osrm\_client.py}, \texttt{lpmc\_inference.py}), y un módulo raíz \texttt{main.py} que monta los routers y configura CORS.

Este diseño tiene una consecuencia práctica relevante: los routers pueden desarrollarse, probarse y documentarse de forma independiente. FastAPI genera documentación OpenAPI automática (\texttt{/docs}) que refleja el estado actual de los endpoints en tiempo real, lo que facilita la depuración y el análisis de la interfaz durante el desarrollo.

\subsection{Router OSRM: consultas multiperfil}

El router \texttt{/api/osrm/routes} acepta un par origen-destino y una lista de perfiles solicitados, lanza las consultas a los tres contenedores OSRM en paralelo mediante \texttt{asyncio.gather} y devuelve los resultados combinados. La URL de cada instancia OSRM se configura mediante variables de entorno (\texttt{OSRM\_DRIVING\_URL}, \texttt{OSRM\_CYCLING\_URL}, \texttt{OSRM\_FOOT\_URL}), lo que permite cambiar los endpoints sin modificar código.

Para cada ruta, el backend decodifica la geometría polyline devuelta por OSRM y la transforma en una lista de coordenadas \texttt{[lat, lon]} lista para consumo directo por Leaflet en el frontend.

\subsection{Router OTP: itinerarios multimodales}

El router \texttt{/api/otp/routes} consulta el endpoint de planificación de OTP, gestiona la selección de itinerario (automática o por índice) y normaliza la respuesta. Para cada \textit{leg} del itinerario devuelve: modo de transporte, distancia, duración, geometría, parada de inicio y fin, nombre de línea, agencia y horarios de salida y llegada formateados en \texttt{HH:MM}.

La geometría de cada tramo se decodifica desde el formato polyline que devuelve OTP y se concatena para formar la traza completa del itinerario, que el frontend puede mostrar como una polilínea continua o como segmentos coloreados por modo.

\subsection{Decisión de entorno Docker}

El backend se ejecuta en un contenedor Docker que monta el repositorio completo en \texttt{/workspace}. Esta configuración permite que el código del backend y los artefactos del modelo LPMC (almacenados en \texttt{lpmc/models/}) sean accesibles desde el mismo sistema de ficheros sin necesidad de copiarlos al contenedor. La función \texttt{\_project\_root()} en el módulo de inferencia resuelve la ruta del modelo usando \texttt{Path(\_\_file\_\_).resolve().parents[4]}, que navega desde el módulo hasta la raíz del repositorio y permite localizar el directorio \texttt{lpmc/models/} de forma independiente del punto de montaje.

% -------------------------------------------------------
\section{Implementación del frontend}
\label{sec:impl_frontend}
% -------------------------------------------------------

\subsection{Tecnologías y estructura}

El frontend está desarrollado con React 18, Vite y TypeScript. Vite actúa como servidor de desarrollo y bundler: proporciona recarga en caliente instantánea durante el desarrollo y genera el bundle de producción con división de código automática. TypeScript añade comprobación estática de tipos, lo que resulta especialmente útil al trabajar con las estructuras de datos complejas que devuelven los endpoints del backend (itinerarios OTP con desglose por tramos, respuestas GTFS con múltiples entidades relacionadas, vectores de probabilidades modales).

La aplicación está estructurada en dos componentes principales: \texttt{App} concentra el estado global y la lógica de negocio, y \texttt{MapView} encapsula la representación cartográfica con React-Leaflet. La separación es deliberada: \texttt{MapView} recibe todos sus datos como \textit{props} y no tiene estado propio, lo que facilita el razonamiento sobre la renderización del mapa.

\subsection{Interfaz cartográfica}

La capa cartográfica usa Leaflet a través de React-Leaflet. El mapa permite al usuario hacer clic para colocar los puntos de origen y destino de forma alternada. Los marcadores de inicio y fin están implementados con iconos personalizados creados con CSS, distintos a los marcadores estándar de Leaflet, para mejorar la identificabilidad visual de ambos puntos.

Las rutas viarias (OSRM) se representan como polilíneas coloreadas según el modo activo: azul para coche, verde para bicicleta, gris para a pie. Los segmentos del itinerario OTP se representan solo cuando el modo activo es «Transporte público», con distinción visual entre tramos a pie (línea discontinua gris) y tramos en autobús (línea continua de color). Este comportamiento evita el solapamiento visual con las rutas viarias cuando el usuario alterna entre modos.

Las paradas del feed GTFS se muestran como marcadores interactivos sobre el mapa; al hacer clic en una parada se abre un popup con el nombre, código y las líneas que pasan por ella. Al seleccionar una línea desde el popup o desde el desplegable del panel, el mapa dibuja el shape de la ruta y el panel lateral muestra el detalle de paradas y los horarios del día seleccionado.

La aplicación incluye cuatro mapas base seleccionables: analítico (CartoDB Light, sin distracciones visuales), color estándar (OpenStreetMap), relieve con curvas de nivel (OpenTopoMap) y vista aérea (Esri World Imagery).

\subsection{Gestión del estado con TanStack Query}

Las peticiones al backend se gestionan con TanStack Query (\texttt{@tanstack/react-query}). El estado de las rutas OSRM y OTP se gestiona con \texttt{useMutation} (petición manual al pulsar el botón «Calcular rutas»), mientras que los datos GTFS estáticos (paradas, lista de rutas) se cargan con \texttt{useQuery} al arrancar la aplicación. TanStack Query proporciona caché automática, estados de carga y error, y reintento con \textit{backoff} sin necesidad de gestión manual.

Las mutaciones de inferencia LPMC (\texttt{/predict}, \texttt{/compare}, \texttt{/debug-features}) se declaran como mutaciones independientes, lo que permite que el usuario invoque predicción simple o comparación de tres modelos de forma independiente sobre el mismo par origen-destino.

\subsection{Perfiles de usuario predefinidos}

Para facilitar la exploración interactiva del simulador sin necesidad de introducir manualmente todos los parámetros sociodemográficos, se implementaron tres perfiles de usuario predefinidos:

\begin{itemize}
  \item \textbf{Commuter}: viaje al trabajo en martes, 8:15, varón de 36 años con carnet de conducir y un coche en el hogar.
  \item \textbf{Estudiante}: viaje a la educación en miércoles, 7:45, mujer de 21 años sin carnet de conducir ni coche en el hogar, con tarifa de transporte público reducida.
  \item \textbf{Familiar}: viaje de ocio en sábado, 11:30, mujer de 44 años con carnet de conducir y dos coches en el hogar.
\end{itemize}

Al activar un perfil se carga el conjunto completo de valores en el formulario de la interfaz. El usuario puede entonces modificar cualquier parámetro individualmente para explorar la sensibilidad de la predicción. La activación del perfil se detecta comparando el estado actual del formulario con los valores predefinidos, por lo que el indicador visual de «perfil activo» desaparece en cuanto el usuario modifica algún campo, evitando confusión sobre qué conjunto de valores se está usando.

% -------------------------------------------------------
\section{Dataset LPMC: descripción y contexto}
\label{sec:impl_lpmc_dataset}
% -------------------------------------------------------

El dataset London Passenger Mode Choice (LPMC) \cite{Hillel2018LPMC,CSLPMC2019} contiene registros de viajes observados en el área metropolitana de Londres, enriquecidos con variables de nivel de servicio calculadas para cada modo de transporte disponible en cada par origen-destino. El dataset fue construido por Hillel et al. como banco de pruebas para la comparación de modelos de elección discreta y de aprendizaje automático en el dominio del transporte, y es la fuente del trabajo de investigación del director del TFM \cite{MartinBaos2023Thesis,MartinBaos2023TRC}.

El archivo original contiene 81.086 registros y 17.616 hogares únicos identificados por \texttt{household\_id}. Los viajes corresponden a los años 2012-2015 y están agrupados en tres oleadas de encuesta (\texttt{survey\_year} 1, 2, 3). El vector de entrada incluye variables de nivel de servicio (tiempos y costes por modo), variables sociodemográficas y contextuales, y la variable objetivo \texttt{travel\_mode} con cuatro categorías: \textit{walk}, \textit{cycle}, \textit{pt} (transporte público) y \textit{drive} (coche).

La Tabla~\ref{tab:lpmc_dataset_stats} resume las estadísticas básicas del dataset.

\begin{table}[htbp]
  \caption{Características generales del dataset LPMC.}
  \label{tab:lpmc_dataset_stats}
  \centering
  \begin{tabular}{lc}
    \toprule
    \textbf{Característica} & \textbf{Valor} \\
    \midrule
    Registros totales        & 81.086 \\
    Hogares únicos           & 17.616 \\
    Modos de viaje           & 4 (walk, cycle, pt, drive) \\
    Oleadas de encuesta      & 3 (survey\_year 1, 2, 3) \\
    Años de viaje cubiertos  & 2012--2015 \\
    Variables originales     & 34 \\
    \bottomrule
  \end{tabular}
\end{table}

El interés metodológico del dataset para este TFM no reside en el caso geográfico concreto (Londres), sino en tres características que lo hacen adecuado como banco de pruebas: dispone de variables de nivel de servicio equivalentes a las que OSRM y OTP pueden calcular para cualquier par origen-destino, incluye variables sociodemográficas similares a las recogidas por encuestas de movilidad estándar, y tiene un tamaño suficiente para entrenar modelos de aprendizaje automático con validación cruzada robusta.

% -------------------------------------------------------
\section{Preprocesado de datos}
\label{sec:impl_preprocesado}
% -------------------------------------------------------

\subsection{Partición temporal train/test}

La separación entre conjunto de entrenamiento y conjunto de test se realiza por oleada de encuesta: las oleadas 1 y 2 forman el conjunto de entrenamiento y la oleada 3 forma el conjunto de test independiente. Esta estrategia de partición temporal evita la fuga de información (\textit{data leakage}) que ocurriría si los viajes del mismo hogar o persona aparecieran simultáneamente en train y en test: al respetar el orden temporal de las oleadas, el modelo no ve en entrenamiento comportamientos observados en la encuesta más reciente.

\subsection{Transformaciones y codificación categórica}

El preprocesado del dataset incluye las siguientes transformaciones:

\begin{itemize}
  \item \textbf{Variable objetivo}: la columna \texttt{travel\_mode} se codifica numéricamente: 0 = \textit{walk}, 1 = \textit{cycle}, 2 = \textit{pt}, 3 = \textit{drive}.
  \item \textbf{Motivo del viaje} (\texttt{purpose}): codificación \textit{one-hot} en cinco columnas (\texttt{purpose\_B}, \texttt{purpose\_HBE}, \texttt{purpose\_HBO}, \texttt{purpose\_HBW}, \texttt{purpose\_NHBO}).
  \item \textbf{Tipo de combustible} (\texttt{fueltype}): los seis valores originales se reagrupan en cuatro categorías para reducir dimensionalidad (\texttt{Petrol\_Car} y \texttt{Petrol\_LGV} → \texttt{Petrol}; \texttt{Diesel\_Car} y \texttt{Diesel\_LGV} → \texttt{Diesel}; \texttt{Hybrid\_Car} → \texttt{Hybrid}; \texttt{Average\_Car} → \texttt{Average}), y se codifican en cuatro columnas \textit{one-hot}.
  \item \textbf{Variables eliminadas}: se descartan \texttt{faretype} (no disponible en tiempo de inferencia), \texttt{trip\_id}, \texttt{person\_n}, \texttt{trip\_n}, \texttt{travel\_year}, \texttt{travel\_month}, \texttt{travel\_date}, \texttt{bus\_scale}, \texttt{dur\_pt\_total}, \texttt{dur\_pt\_int\_total}, \texttt{cost\_driving\_fuel}, \texttt{cost\_driving\_con\_charge} y \texttt{driving\_traffic\_percent}. La mayoría corresponde a identificadores o a variables derivadas redundantes; \texttt{faretype} se descarta porque no forma parte de la información que el usuario puede proporcionar al simulador.
  \item \textbf{Identificador de hogar}: \texttt{household\_id} se conserva en el fichero de entrenamiento pero \textbf{nunca entra como variable predictora}. Su único papel es definir los grupos en la validación cruzada agrupada.
\end{itemize}

Tras el preprocesado, el vector de características tiene 28 columnas: 10 variables de ruta, 8 variables sociodemográficas y contextuales directas, 5 columnas de propósito y 4 columnas de tipo de combustible (más \texttt{household\_id} que se excluye antes del entrenamiento), lo que resulta en 27 features efectivas.

\subsection{Escalado parcial con StandardScaler}

El escalado se aplica solo a las 16 variables continuas con distribución asimétrica o rangos muy dispares (Tabla~\ref{tab:scaled_features}). Las variables binarias (\texttt{female}, \texttt{driving\_license}) y los vectores \textit{one-hot} no se escalan, ya que ya están en el rango $[0, 1]$ y el escalado no aportaría información adicional sino ruido.

\begin{table}[htbp]
  \caption{Variables continuas sometidas a escalado con StandardScaler.}
  \label{tab:scaled_features}
  \centering
  \begin{tabular}{ll}
    \toprule
    \textbf{Variable} & \textbf{Descripción} \\
    \midrule
    \texttt{day\_of\_week}           & Día de la semana (1--7) \\
    \texttt{start\_time\_linear}     & Hora de inicio del viaje (decimal) \\
    \texttt{age}                     & Edad del viajero \\
    \texttt{car\_ownership}          & Número de coches en el hogar \\
    \texttt{distance}                & Distancia en coche (m) \\
    \texttt{dur\_walking}            & Tiempo a pie hasta destino (h) \\
    \texttt{dur\_cycling}            & Tiempo en bicicleta (h) \\
    \texttt{dur\_driving}            & Tiempo de conducción (h) \\
    \texttt{dur\_pt\_access}         & Tiempo de acceso peatonal al transporte público (h) \\
    \texttt{dur\_pt\_bus}            & Tiempo en autobús (h) \\
    \texttt{dur\_pt\_rail}           & Tiempo en ferroviario (h) \\
    \texttt{dur\_pt\_int\_waiting}   & Tiempo de espera en transbordos (h) \\
    \texttt{dur\_pt\_int\_walking}   & Desplazamiento a pie en transbordos (h) \\
    \texttt{pt\_n\_interchanges}     & Número de transbordos \\
    \texttt{cost\_transit}           & Coste del billete de transporte público (£) \\
    \texttt{cost\_driving\_total}    & Coste total estimado del viaje en coche (£) \\
    \bottomrule
  \end{tabular}
\end{table}

El escalador se ajusta (\texttt{fit}) exclusivamente sobre el subconjunto de entrenamiento del fold o del dataset final, y se aplica (\texttt{transform}) tanto al conjunto de entrenamiento como al de validación o test. Este protocolo evita que la información estadística del conjunto de test contamine el ajuste del escalador. El escalador ajustado sobre el conjunto de entrenamiento completo se serializa junto con el modelo y se carga en el pipeline de inferencia del backend para garantizar que la normalización en producción es idéntica a la usada en el entrenamiento.

% -------------------------------------------------------
\section{Entrenamiento y validación de los modelos}
\label{sec:impl_entrenamiento}
% -------------------------------------------------------

\subsection{Validación cruzada agrupada por hogar}

Un aspecto metodológico central del proceso de entrenamiento es el uso de \texttt{GroupKFold} con \texttt{n\_splits=5} agrupando los viajes por \texttt{household\_id}. Esta estrategia garantiza que todos los viajes de un mismo hogar quedan asignados al mismo fold, de modo que nunca aparecen simultáneamente en el fold de entrenamiento y en el fold de validación dentro de un mismo split.

La motivación es evitar optimismo artificial en las métricas de validación cruzada. Si los viajes del mismo hogar se distribuyen aleatoriamente entre folds, el modelo puede aprender patrones idiosincrásicos del hogar y reflejarlos en la validación del mismo fold, lo que produce estimaciones de rendimiento irrealmente altas que no se sostienen sobre datos de hogares no vistos. Con \texttt{GroupKFold}, las métricas de la CV reflejan el comportamiento esperado sobre hogares completamente nuevos.

Este mismo razonamiento aplica al escalado: el \texttt{StandardScaler} se ajusta solo sobre los datos de entrenamiento del fold, y se transforma el fold de validación con los parámetros aprendidos en el fold de entrenamiento. Ajustar el escalador sobre todo el fold (incluyendo validación) sería otra forma de fuga de información, más sutil pero igualmente distorsionadora.

\subsection{XGBoost: configuración e hiperparámetros}

Los hiperparámetros de XGBoost se obtuvieron de la investigación previa del director del TFM, que realizó una búsqueda exhaustiva sobre el mismo dataset LPMC con 1.000 evaluaciones de validación cruzada de 5-folds \cite{MartinBaos2023Thesis,MartinBaos2023TRC}. Los parámetros resultantes se almacenan en \texttt{lpmc/artifacts/lpmc\_xgb\_best\_params.json} y se cargaron directamente en el script de entrenamiento. La Tabla~\ref{tab:xgb_params} recoge los principales hiperparámetros.

\begin{table}[htbp]
  \caption{Hiperparámetros del modelo XGBoost.}
  \label{tab:xgb_params}
  \centering
  \begin{tabular}{ll}
    \toprule
    \textbf{Hiperparámetro} & \textbf{Valor} \\
    \midrule
    \texttt{n\_estimators}      & 4.809 \\
    \texttt{learning\_rate}     & 0,01 \\
    \texttt{max\_depth}         & 4 \\
    \texttt{gamma}              & 3,93 \\
    \texttt{min\_child\_weight} & 35 \\
    \texttt{subsample}          & 0,767 \\
    \texttt{colsample\_bytree}  & 0,934 \\
    \texttt{colsample\_bylevel} & 0,651 \\
    \texttt{reg\_alpha}         & 0,048 \\
    \texttt{reg\_lambda}        & 0,037 \\
    \texttt{objective}          & \texttt{multi:softprob} \\
    \bottomrule
  \end{tabular}
\end{table}

El número de estimadores elevado (4.809) combinado con una tasa de aprendizaje baja (0,01) es un patrón habitual en XGBoost optimizado: muchos árboles poco profundos con poca tasa de aprendizaje evitan el sobreajuste manteniendo alta la capacidad del ensemble.

\subsection{Random Forest: configuración}

Random Forest se entrenó con 500 árboles sin límite de profundidad (\texttt{max\_depth=None}), con \texttt{max\_features=\textquotesingle{}sqrt\textquotesingle{}} para la selección aleatoria de variables en cada split, y restricciones mínimas de tamaño de nodo (\texttt{min\_samples\_split=5}, \texttt{min\_samples\_leaf=2}) para controlar el sobreajuste sin limitar el crecimiento de los árboles. La semilla aleatoria es 481516 para reproducibilidad.

Random Forest es inherentemente invariante a la escala de las variables de entrada porque sus splits se basan en comparaciones de orden, no en magnitudes absolutas. Aún así, el escalado se aplica de forma idéntica a los tres modelos para mantener la consistencia del pipeline y facilitar la depuración del vector de características en producción.

\subsection{Red neuronal profunda (DNN)}

La DNN está implementada con PyTorch y sigue la arquitectura de la Tabla~\ref{tab:dnn_arch}. El patrón \texttt{Linear~→~BatchNorm~→~ReLU} (BatchNorm después de la transformación lineal y antes de la activación) normaliza las activaciones de cada capa y estabiliza el entrenamiento, especialmente con los rangos de magnitud heterogéneos que tienen las variables del dataset LPMC.

\begin{table}[htbp]
  \caption{Arquitectura de la red neuronal profunda.}
  \label{tab:dnn_arch}
  \centering
  \begin{tabular}{lll}
    \toprule
    \textbf{Capa} & \textbf{Dimensiones} & \textbf{Detalles} \\
    \midrule
    Linear        & $n\_\text{features} \to 128$ & \\
    BatchNorm1d   & 128                  & \\
    ReLU          &                      & \\
    Dropout       &                      & $p = 0{,}3$ \\
    Linear        & $128 \to 64$         & \\
    BatchNorm1d   & 64                   & \\
    ReLU          &                      & \\
    Dropout       &                      & $p = 0{,}2$ \\
    Linear        & $64 \to 32$          & \\
    BatchNorm1d   & 32                   & \\
    ReLU          &                      & \\
    Linear        & $32 \to 4$           & logits de salida \\
    \bottomrule
  \end{tabular}
\end{table}

El entrenamiento usa el optimizador Adam con \texttt{lr=0,001} y \texttt{weight\_decay=0,001} para regularización L2. La función de pérdida es \texttt{CrossEntropyLoss} con \texttt{label\_smoothing=0,1}: el suavizado de etiquetas reemplaza los objetivos duros (0 o 1) por distribuciones suavizadas ($1 - \varepsilon$ para la clase correcta y $\varepsilon / (K-1)$ para las incorrectas, con $\varepsilon = 0{,}1$ y $K = 4$), lo que evita que el modelo tienda a maximizar logits indefinidamente y mejora la calibración de las probabilidades de salida.

El gradiente se recorta a norma máxima 1 en cada paso (\texttt{clip\_grad\_norm\_}), lo que estabiliza el entrenamiento en las primeras épocas cuando los pesos aún no han convergido. El scheduler \texttt{ReduceLROnPlateau} reduce la tasa de aprendizaje a la mitad cuando la pérdida de validación no mejora durante 5 épocas consecutivas. El entrenamiento se detiene anticipadamente (\textit{early stopping}) si la pérdida de validación no mejora durante 10 épocas, restaurando el estado del modelo con menor pérdida de validación observada. En el entrenamiento del modelo definitivo, el proceso convergió en 60 épocas (de las 100 configuradas como máximo).

El modelo final se serializa en dos partes: el \texttt{state\_dict} de PyTorch en un archivo \texttt{.pt} y un bundle \texttt{.joblib} que almacena la ruta al archivo \texttt{.pt}, el número de features y los nombres de las columnas. Esta separación permite cargar la arquitectura desde el código del backend sin necesidad de serializar el objeto PyTorch completo, lo que facilita la portabilidad entre entornos.

% -------------------------------------------------------
\section{Pipeline de inferencia en producción}
\label{sec:impl_pipeline}
% -------------------------------------------------------

\subsection{Construcción concurrente del vector de características}

Al recibir una petición de inferencia, el backend lanza en paralelo cuatro consultas asíncronas mediante \texttt{asyncio.gather}: una a cada una de las tres instancias OSRM (conducción, bicicleta, a pie) y una a OpenTripPlanner. La paralelización reduce la latencia total de la petición al tiempo de la consulta más lenta, en lugar de la suma de todas.

La función \texttt{\_build\_route\_features} extrae del resultado de OSRM y OTP las diez variables de ruta descritas en la Tabla~\ref{tab:lpmc_features} del capítulo~\ref{ch:arquitectura}. La función \texttt{\_build\_feature\_frame} combina las variables de ruta con el perfil sociodemográfico recibido en la petición, aplica la codificación \textit{one-hot} a las variables categóricas y construye el array de entrada al modelo con las columnas en el mismo orden que en el entrenamiento.

\subsection{Corrección de unidades: bug y resolución}
\label{subsec:bug_unidades}

Durante el desarrollo del pipeline de inferencia se detectó un error crítico de unidades: OSRM y OTP devuelven duraciones en segundos, pero el dataset LPMC almacena todas las duraciones de viaje en horas. El error se manifestó inicialmente de forma silenciosa en los modelos de árbol (XGBoost y Random Forest), que son invariantes a la escala por construcción y producen predicciones estadísticamente plausibles aunque incorrectas. La DNN, sin embargo, falló visiblemente: el escalador se había ajustado con valores de orden 1 (horas) y recibía en producción valores de orden 3600 (segundos), lo que producía inputs a más de 700 desviaciones típicas de la media del escalador. En esas condiciones, el softmax de salida colapsaba con probabilidad 1,0 para un solo modo.

La corrección se realizó en la función \texttt{\_build\_route\_features} aplicando la conversión \texttt{s2h = 1.0 / 3600.0} a todas las duraciones antes de construir el vector:

\begin{equation}
  d_{\text{hours}} = \frac{d_{\text{seconds}}}{3600}
\end{equation}

Con esta corrección, los tres modelos reciben el mismo vector de características y producen predicciones correctas. La detección del bug fue un resultado positivo del uso de tres modelos: la DNN actuó como señal de alerta de un error que los modelos de árbol enmascaraban.

\subsection{Carga perezosa y caché de artefactos}

Los artefactos de modelo (bundle joblib + estado PyTorch) no se cargan al arrancar el backend sino en la primera petición que los requiere. El módulo de inferencia mantiene un diccionario \texttt{\_ARTIFACTS\_CACHE} indexado por variante (\texttt{xgb}, \texttt{rf}, \texttt{dnn}): la primera petición para cada variante carga y deserializa los artefactos; las peticiones siguientes los reutilizan desde memoria.

Para la DNN, la carga es en dos fases: el bundle joblib almacena la ruta al archivo \texttt{.pt} (no el tensor en sí); al crear el wrapper \texttt{TorchModalWrapper}, el estado del modelo se carga desde disco solo cuando se llama a \texttt{predict\_proba} por primera vez. Esto permite que el bundle sea portable entre sistemas operativos: si la ruta almacenada en el bundle es una ruta absoluta de Windows que no existe en el contenedor Linux, el código normaliza el separador de ruta antes de extraer el nombre del fichero (sustituyendo \texttt{\textbackslash{}} por \texttt{/}) y lo busca en el directorio estándar de modelos del repositorio.

\subsection{Modo de comparación de tres modelos}

El endpoint \texttt{/api/lpmc/compare} ejecuta la inferencia con los tres modelos sobre el mismo vector de características y devuelve las distribuciones de probabilidad de los tres en una sola respuesta. La implementación reutiliza todas las funciones del pipeline de predicción individual: construye el vector de características una sola vez, y luego invoca \texttt{\_predict(x, feature\_names, variant)} para cada variante con el modelo y escalador correspondientes.

El frontend muestra los resultados en una tabla comparativa que identifica visualmente el modo predicho por cada modelo. Esto permite al usuario evaluar la consistencia entre modelos para el mismo viaje y perfil, lo que es útil tanto para la exploración interactiva como para detectar casos en los que los modelos discrepan.

% -------------------------------------------------------
\section{Resultados de los modelos}
\label{sec:impl_resultados}
% -------------------------------------------------------

\subsection{Métricas de evaluación}

Los modelos se evalúan con dos métricas:

\begin{itemize}
  \item \textbf{Accuracy}: fracción de predicciones correctas sobre el total. Es la métrica más intuitiva pero puede ser engañosa con distribuciones de clase desequilibradas.
  \item \textbf{GMPCA} (\textit{Geometric Mean Probability of Correct Alternative}): media geométrica de la probabilidad asignada a la alternativa correcta por el modelo. Equivale a \(\exp(-H_{\text{CE}})\), donde $H_{\text{CE}}$ es la entropía cruzada del modelo. Esta métrica penaliza predicciones con alta confianza en la clase incorrecta más severamente que el accuracy, y es estándar en la evaluación de modelos de elección discreta \cite{MartinBaos2023Thesis,MartinBaos2023TRC}.
\end{itemize}

Las métricas de entrenamiento se estiman mediante la media de los 5 folds de la validación cruzada agrupada descrita anteriormente.

\subsection{Resultados comparativos}

Las tablas~\ref{tab:lpmc_train_results} y~\ref{tab:lpmc_test_results} presentan los resultados de los tres modelos en entrenamiento (validación cruzada) y en el conjunto de test independiente, respectivamente.

\begin{table}[htbp]
  \caption{Resultados en entrenamiento con validación cruzada de 5-folds agrupada por hogar (dataset LPMC).}
  \label{tab:lpmc_train_results}
  \centering
  \begin{tabular}{lcc}
    \toprule
    \textbf{Modelo} & \textbf{Accuracy (\%)} & \textbf{GMPCA (\%)} \\
    \midrule
    Random Forest & 74,87 & 51,18 \\
    XGBoost       & 75,48 & 52,54 \\
    DNN           & 75,21 & 51,22 \\
    \bottomrule
  \end{tabular}
\end{table}

\begin{table}[htbp]
  \caption{Resultados en el conjunto de test independiente (dataset LPMC).}
  \label{tab:lpmc_test_results}
  \centering
  \begin{tabular}{lcc}
    \toprule
    \textbf{Modelo} & \textbf{Accuracy (\%)} & \textbf{GMPCA (\%)} \\
    \midrule
    Random Forest & 74,10 & 50,33 \\
    XGBoost       & 74,41 & 51,63 \\
    DNN           & 74,27 & 50,37 \\
    \bottomrule
  \end{tabular}
\end{table}

Los tres modelos obtienen resultados muy similares tanto en accuracy como en GMPCA. XGBoost supera ligeramente a Random Forest y a la DNN en ambas métricas y en ambos conjuntos, siendo el más adecuado como modelo activo por defecto. La DNN, a pesar de ser el modelo más complejo, no supera a XGBoost en este dataset, lo que es coherente con la literatura que documenta que los métodos de boosting suelen competir con redes neuronales en datos tabulares estructurados de tamaño moderado \cite{Grinsztajn2022TreesVsNNs}. La diferencia entre las métricas de validación cruzada y las de test es pequeña (inferior al 1\%), lo que indica que los modelos generalizan bien y el protocolo de validación agrupada por hogar es adecuado para estimar el rendimiento en producción.

\subsection{Interpretación de los perfiles de usuario}

La tabla comparativa de tres modelos en el frontend permite explorar de forma cualitativa la sensibilidad de los modelos a los parámetros del perfil de usuario. Algunos comportamientos observados al variar los perfiles predefinidos:

\begin{itemize}
  \item \textbf{Commuter} con coche disponible: los tres modelos coinciden en predecir \textit{drive} para trayectos de distancia media-alta cuando el coste del coche es moderado.
  \item \textbf{Estudiante} sin carnet: la ausencia de \texttt{driving\_license} y \texttt{car\_ownership=0} penaliza fuertemente la opción \textit{drive}. Para trayectos cortos y con parada de autobús próxima, la predicción converge a \textit{pt} o \textit{walk} según los tiempos de acceso al transporte público.
  \item \textbf{Familiar} en sábado: el día de la semana influye en la disponibilidad de transporte público, lo que puede desplazar la predicción hacia \textit{drive} para itinerarios con frecuencia de bus reducida en fin de semana.
\end{itemize}

Estos comportamientos son cualitativamente coherentes con la lógica de elección modal descrita en la literatura: la disponibilidad de vehículo privado, el tiempo de acceso al transporte público y el coste relativo son los principales determinantes de la elección \cite{Train2009,MartinBaos2023Thesis}.

\subsection{Limitación: dominio de entrenamiento versus área de aplicación}

Los modelos se entrenaron con datos de Londres (2012-2015) y se aplican a viajes en Toledo (2025). Esta discrepancia de dominio es una limitación explícita del simulador: los valores absolutos de las probabilidades y el modo predicho deben interpretarse como indicativos de la tendencia relativa de elección, no como estimaciones cuantitativas precisas para el contexto toledano. El interés del simulador en este TFM es demostrar la viabilidad de integrar modelos de elección modal con enrutado multimodal en un entorno web interactivo, no producir predicciones calibradas para Toledo.

El uso del dataset LPMC está motivado por la disponibilidad de datos con variables de nivel de servicio compatibles con las que proporcionan OSRM y OTP, por su empleo previo en la investigación del director del TFM \cite{MartinBaos2023Thesis}, y por la ausencia de un dataset equivalente con datos reales de Toledo o Castilla-La Mancha.
